\chapter{Arquitectura del Sistema}

\section{Monorepo con Turborepo}

Vetrinaria utiliza una arquitectura de monorepo gestionada con \textbf{pnpm
workspaces} y orquestada por \textbf{Turborepo}. Esta elecci\'on permite
compartir configuraciones, tipos y utilidades entre m\'ultiples paquetes
mientras se mantiene un \'unico repositorio.

\subsection{Estructura de Directorios}

\begin{lstlisting}[style=bash, caption=Estructura completa del monorepo]
vetrinaria/
+-- apps/
|   +-- web/                    # Next.js 15 App Router
|       +-- src/
|           +-- app/            # Rutas y layouts
|           |   +-- (dashboard)/# Layout con sidebar + navbar
|           |   +-- auth/       # Login y registro
|           |   +-- api/        # API routes
|           +-- components/     # Componentes React
|           |   +-- ui/         # shadcn/ui primitives
|           +-- hooks/          # Custom hooks
|           +-- lib/            # Utilidades y config
|               +-- actions/    # Server Actions
+-- packages/
|   +-- db/                     # Drizzle ORM + esquemas + seed
|   +-- shared/                 # Tipos y constantes compartidos
|   +-- config/                 # ESLint, tsconfig base
|   +-- ui/                     # Placeholder para componentes UI
+-- docker-compose.yml          # Servicios PostgreSQL, MinIO, Redis
+-- turbo.json                  # Configuracion de Turborepo
\end{lstlisting}

\subsection{Configuraci\'on de Turborepo}

El archivo \texttt{turbo.json} define las tareas del pipeline de build:

\begin{lstlisting}[style=json, caption=turbo.json]
{
  "tasks": {
    "dev":        { "cache": false, "persistent": true },
    "build":      { "dependsOn": ["^build"],
                    "outputs": [".next/**", "dist/**"] },
    "lint":       { "dependsOn": ["^build"] },
    "typecheck":  { "dependsOn": ["^build"] },
    "generate":   {},
    "migrate":    {},
    "push":       {},
    "studio":     { "cache": false, "persistent": true },
    "seed":       { "dependsOn": ["^build"] },
    "test":       {}
  }
}
\end{lstlisting}

Las caracter\'{\i}sticas clave de Turborepo en este proyecto son:

\begin{itemize}[leftmargin=*]
    \item \textbf{Caching inteligente:} Las tareas \texttt{build}, \texttt{lint}
          y \texttt{typecheck} se cachean, evitando recomputaciones innecesarias.
    \item \textbf{Ejecuci\'on paralela:} Tareas independientes se ejecutan
          simult\'aneamente.
    \item \textbf{Pipeline de dependencias:} \texttt{build} depende de
          \texttt{\^build} (build de dependencias primero).
    \item \textbf{Persistencia:} \texttt{dev} y \texttt{studio} son tareas
          persistentes que no se cachean.
\end{itemize}

\section{Workspaces de pnpm}

El archivo \texttt{pnpm-workspace.yaml} declara los workspaces:

\begin{lstlisting}[style=yaml, caption=pnpm-workspace.yaml]
packages:
  - 'apps/*'
  - 'packages/*'
\end{lstlisting}

Las dependencias compartidas se declaran en el \texttt{package.json} ra\'{\i}z:

\begin{lstlisting}[style=json, caption=package.json raiz (scripts)]
{
  "scripts": {
    "dev":       "turbo dev",
    "build":     "turbo build",
    "lint":      "turbo lint",
    "typecheck": "turbo typecheck",
    "format":    "prettier --write .",
    "test":      "turbo test",
    "db:generate": "pnpm --filter @vetrinaria/db generate",
    "db:migrate":  "pnpm --filter @vetrinaria/db migrate",
    "db:push":     "pnpm --filter @vetrinaria/db push",
    "db:studio":   "pnpm --filter @vetrinaria/db studio",
    "db:seed":     "pnpm --filter @vetrinaria/db seed"
  }
}
\end{lstlisting}

\section{Paquetes del Monorepo}

\subsection{@vetrinaria/web (apps/web)}

Aplicaci\'on principal Next.js 15 con App Router. Contiene:

\begin{itemize}[leftmargin=*]
    \item \textbf{P\'aginas:} Dashboard, autenticaci\'on y m\'odulos de negocio.
    \item \textbf{API Routes:} Endpoints REST para el cliente.
    \item \textbf{Server Actions:} Mutaciones de datos del lado del servidor.
    \item \textbf{Componentes:} UI primitives (shadcn/ui) y componentes de negocio.
    \item \textbf{Configuraci\'on:} NextAuth, variables de entorno, middleware.
\end{itemize}

La configuraci\'on de Next.js transpila los paquetes del workspace:

\begin{lstlisting}[style=typescript, caption=next.config.ts]
const nextConfig: NextConfig = {
  transpilePackages: ['@vetrinaria/db', '@vetrinaria/shared'],
  experimental: {
    optimizePackageImports: ['lucide-react', '@radix-ui/react-*'],
    turbo: {
      resolveExtensions: ['.ts', '.tsx', '.js', '.jsx', '.json'],
    },
  },
  images: {
    formats: ['image/avif', 'image/webp'],
  },
};
\end{lstlisting}

\subsection{@vetrinaria/db (packages/db)}

Paquete de base de datos que contiene:

\begin{itemize}[leftmargin=*]
    \item Definiciones de esquemas Drizzle ORM (14 tablas).
    \item Configuraci\'on de conexi\'on PostgreSQL.
    \item Migraciones generadas por Drizzle Kit.
    \item Script de seed con datos de prueba realistas.
    \item Relaciones entre tablas.
\end{itemize}

Conexi\'on a base de datos con detecci\'on autom\'atica de SSL:

\begin{lstlisting}[style=typescript, caption=Conexion a PostgreSQL con postgres.js]
const connectionString = cleanConnectionString(process.env.DATABASE_URL!);

const needsSsl = connectionString.includes('neon.tech')
                 || process.env.NODE_ENV === 'production';

const client = postgres(connectionString, {
  prepare: false,
  max: 1,
  idle_timeout: 20,
  connect_timeout: 30,
  ...(needsSsl ? { ssl: { rejectUnauthorized: false } } : {}),
});

export const db = drizzle(client, { schema, logger: false });
export { schema };
\end{lstlisting}

\subsection{@vetrinaria/shared (packages/shared)}

Tipos y constantes compartidas entre todos los paquetes:

\begin{itemize}[leftmargin=*]
    \item \textbf{Tipos:} \texttt{UserRole}, \texttt{AppointmentStatus},
          \texttt{AppointmentType}, \texttt{Species}, \texttt{Gender},
          \texttt{ProductType}, \texttt{ToothStatus}.
    \item \textbf{Constantes:} Etiquetas en espa\~nol para todos los enums,
          colores asociados a estados dentales, numeraci\'on Triadan para
          dientes de perros (42) y gatos (30).
\end{itemize}

\subsection{@vetrinaria/config (packages/config)}

Configuraciones base compartidas:

\begin{itemize}[leftmargin=*]
    \item \texttt{tsconfig.base.json}: ES2022, strict mode, bundler resolution.
    \item \texttt{tsconfig.nextjs.json}: Extiende base con jsx:preserve.
    \item \texttt{eslint-base.js}: Reglas ESLint con @typescript-eslint.
\end{itemize}

\subsection{@vetrinaria/ui (packages/ui)}

Paquete placeholder para futuros componentes UI compartidos.

\section{Flujo de Datos}

El sistema sigue una arquitectura de tres capas:

\begin{enumerate}[leftmargin=*]
    \item \textbf{Capa de Presentaci\'on:} Componentes React (Server y Client
          Components) que renderizan la interfaz.
    \item \textbf{Capa de L\'ogica:} Server Actions (mutaciones) y Server
          Components (lecturas) que ejecutan la l\'ogica de negocio.
    \item \textbf{Capa de Datos:} Drizzle ORM sobre PostgreSQL, con esquemas
          definidos en \texttt{packages/db}.
\end{enumerate}

\subsection{Server Actions vs API Routes}

El proyecto utiliza dos patrones para la comunicaci\'on servidor-cliente:

\begin{itemize}[leftmargin=*]
    \item \textbf{Server Actions} (\texttt{'use server'}): Para mutaciones de
          datos desde formularios y componentes del lado del cliente. Ofrecen
          revalidaci\'on autom\'atica de rutas via \texttt{revalidatePath()}.
    \item \textbf{API Routes} (\texttt{app/api/}): Para endpoints REST
          consumidos por JavaScript del lado del cliente (fetch). Se usan
          principalmente para cargar datos en selectores dependientes.
\end{itemize}

\section{TypeScript y Type Safety}

El proyecto utiliza TypeScript 5.8 en modo estricto con configuraci\'on base:

\begin{lstlisting}[style=json, caption=tsconfig.base.json (extracto)]
{
  "compilerOptions": {
    "target": "ES2022",
    "module": "ES2022",
    "moduleResolution": "bundler",
    "strict": true,
    "esModuleInterop": true,
    "skipLibCheck": true,
    "declaration": true,
    "declarationMap": true,
    "sourceMap": true
  }
}
\end{lstlisting}

El tipado de extremo a extremo se logra mediante:

\begin{itemize}[leftmargin=*]
    \item \textbf{Zod:} Validaci\'on de esquemas en tiempo de ejecuci\'on
          para formularios y variables de entorno.
    \item \textbf{Drizzle ORM:} Queries tipadas que reflejan el esquema de
          base de datos.
    \item \textbf{TypeScript inferido:} Tipos inferidos desde Zod con
          \texttt{z.infer} y desde Drizzle con tipos de las tablas.
    \item \textbf{@t3-oss/env-nextjs:} Validaci\'on de variables de entorno
          con esquemas Zod.
\end{itemize}
